%% PLANTILLA PARA ARTICULOS
\documentclass[10pt,twoside,spanish,a4paper]{article}
%%%%%%%%%%%%%%%%%%%%%%%%%%%%% no tocar
% Alfredo Cuesta. 01/01/2013. acuestai@ucm.es

%si se quiere usar ARIAL hay que poner estas 2 lineas:
%\usepackage[scaled]{uarial} %DESCOMENTAR
\renewcommand*\familydefault{\sfdefault}
%

\usepackage[compact]{titlesec}
 \titlespacing\section{0pt}{12pt plus 0pt minus 0pt}{0pt plus 0pt minus 0pt}
 \titlespacing\subsection{0pt}{5pt plus 0pt minus 0pt}{0pt plus 0pt minus 0pt}
 \titlespacing\subsubsection{0pt}{0pt plus 0pt minus 0pt}{0pt plus 0pt minus 0pt}

\usepackage[utf8]{inputenc}
\usepackage{
amssymb,
graphicx,
color,
url
%wrapfig %Inclusión de gr·ficos al lado de texto
}
\definecolor{olivegreen}{rgb}{0.0,0.5,0.0}

\usepackage{listings}
\lstset{language=Matlab,%
    basicstyle=\footnotesize,
    keywordstyle=\color{blue},%
    commentstyle=\color{olivegreen}\textit,%
    emph=[1]{ans},emphstyle=[1]\color{black}, %some words to emphasise
    emph=[2]{resample,timeseries,syms},emphstyle=[2]\color{blue}, %some words to emphasise
}

\linespread{1.3}%separacion de lineas
\parindent=0.0in%indentancion de parrafo
\parskip=6pt%separacion entre el parrafo siguiente
\setlength{\topmargin}{-0.5in}
\setlength{\oddsidemargin}{0in}
\setlength{\evensidemargin}{0in}
\setlength{\textwidth}{16cm}%{4.72in}
\setlength{\textheight}{22cm}%{7.08in}

\usepackage{babel}
\addto\shorthandsspanish{\spanishdeactivate{~<>}}
% correct bad hyphenation here
\hyphenation{nu-es-tro co-no-ci-mien-to par-ti-cu-lar ca-rac-te-ris-ti-cas ni-vel se-gui-mien-to que-re-mos ca-ta-lo-go}
\date{}


% --- customization ---

\renewcommand\labelitemi{$\bullet$}
\let\olditemize\itemize
\renewcommand{\itemize}{
  \olditemize
  \vspace{-0pt}
  \setlength{\itemsep}{6pt}
  \setlength{\parskip}{0pt}
  \setlength{\parsep}{0pt}
  %\setlength{\topsep}{-15pt}
}

\let\oldenumerate\enumerate
\renewcommand{\enumerate}{
  \oldenumerate
  \vspace{-2pt}
  \setlength{\itemsep}{1pt}
  \setlength{\parskip}{0pt}
  \setlength{\parsep}{0pt}
  %\setlength{\topsep}{-15pt}
}

\usepackage{enumitem}
\usepackage{verbatim}
\usepackage{fancyhdr}
\usepackage{multirow}
\usepackage{colortbl}
\usepackage{verbatim}
\usepackage{amsmath}

\begin{document}

\renewcommand{\tablename}{Tabla}
\renewcommand{\figurename}{Figura}
\title{
  \vspace{-0.2cm}
	\begin{center}
		\large
		\begin{tabular}{rrl}
		\multirow{4}{*}{\includegraphics[width=1.7cm]{UCM.eps} \hspace{1cm}}&Grado: & Ingeniería Electrónica de Comunicaciones\\
		&Asignatura: & Procesamiento de Señales \\
		&Curso:      & 2025/2026 \\
		\end{tabular}
	\end{center}
	\rule{\textwidth}{0.4pt}
	%
	\vspace{-1.5cm}
	\begin{center}
		\color{blue} \Huge
		\begin{tabular}{c}
	        \textbf{Práctica 7}\\
		\textbf{Variables aleatorias}\\
		\end{tabular}
	\end{center}
	\vspace{-0.5cm}
	\rule{\textwidth}{0.4pt}
}
\maketitle
 \thispagestyle{empty}
\pagestyle{fancy}

\lfoot{\includegraphics[width=0.8cm]{UCM.eps}}
\cfoot{\vspace{-0.6cm} Ingeniería Electrónica de Comunicaciones }
\rfoot{\vspace{-0.6cm} \thepage}
\lhead{Práctica}
\rhead{Práctica 7: Variables aleatorias}
\renewcommand{\headrule}{ \vspace{-0.2cm} \rule{\textwidth}{0.4pt}}
\renewcommand{\footrule}{\vspace{0.5cm} \rule{\textwidth}{0.4pt}\\ \vspace{-0.3cm} \\}

%**************************************************************<- para indicar pregunta de la pr·ctica
\newcounter{ctPreg}
\setcounter{ctPreg}{1}
\newcommand{\itemP}{ \item[{\color{blue}\fbox{\textbf{Tarea \arabic{ctPreg}:}}}] \addtocounter{ctPreg}{1}}  %
\newcommand{\itemB}{ \item[{\color{blue}$\bullet$} ]    }
\newcommand{\itemR}{ \item[{\color{blue}$\triangleright$} ]    }       %
%**************************************************************

\vspace{-8em}

\section*{Objetivos}

\begin{itemize}
	\item Aprender a caracterizar el comportamiento de diferentes conjuntos de variables aleatorias.
\end{itemize}

\section*{Generación y análisis de conjuntos de variables aleatorias}
\begin{itemize}
	\itemP \textbf{Muestreo de diferentes conjuntos de variables aleatorias.}

	Generar 100000 muestras de cada conjunto de variables aleatorias a partir de las siguientes funciones de densidad de probabilidad.

	\begin{itemize}
		\itemR Caso 1: $p_{XYZ}(x,y,z)=\mathcal{N}_{(\mu,\Sigma)}(x,y,z)$\\[2ex]
		$\quad\mu=\left[\begin{matrix}2\\8\\-2\end{matrix}\right],~\Sigma=\left[\begin{matrix}2 & 0.5 & 1\\0.5 & 3 & 2\\1 & 2 & 2\end{matrix}\right]$

		\itemR Caso 2: $p_{XYZ}(x,y,z)=\mathcal{N}_{(\mu,\Sigma)}(x,y)\mathcal{U}_{[2,5]}(z)$\\[2ex]
		$\mu=\left[\begin{matrix}2\\8\end{matrix}\right]$
		$\Sigma=\left[\begin{matrix}2 & 3\\3 & 6\end{matrix}\right]$
	\end{itemize}

	Para comprobar que las muestras generadas siguen la función de probabilidad propuesta:
	\begin{enumerate}[label=\alph*)]
		\item Representa gráficamente los valores de las muestras. 
		\item Representa en histogramas 1D (\color{blue}\texttt{histogram}\color{black}), las variables probabilisticamente independientes.
		\item Representa en histogramas 2D (\color{blue}\texttt{histogram2}\color{black}), los pares de variables dependientes.
		\item Compara los histogramas con la función de probabilidad original.
	\end{enumerate}


	\hspace{0.2cm}\color{blue}Ayuda: \color{black} \color{black} \'{U}nicamente se pueden usar los generadores aleatorios de Matlab  \color{blue}randi\color{black}, \color{blue}randn \color{black} y \color{blue}rand\color{black}.  Por lo tanto, se recomienda utilizar el método de muestreo más sencillo que, utilizando dichas funciones, genere las muestras propuestas.

	\itemP  \textbf{Caracterización de conjuntos muestrales desconocidos}

	El objetivo de esta tarea es caracterizar conjuntos muestrales de 2 variables aleatorias de origen desconocido a partir de su histograma 2D. Los 2 conjuntos muestrales que hay que analizar se generarán utilizando el DNI del alumno y la función \color{blue}\texttt{DatosTarea2}\color{black}.
	\begin{enumerate}[label=\alph*)]
		\item Analiza el histograma e intenta identificar la forma de la función de probabilidad original.
		\item Calcula los estadísticos habituales de cada uno de los conjuntos e indica en que casos parametrizan correctamente la función de probabilidad.
		\item Indica en qué casos los estadísticos permiten determinar alguna característica de la función de probabilidad.
	\end{enumerate}
\end{itemize}

\end{document}
